\documentclass[aspectratio=169,10pt,t]{beamer}
\usepackage{ac1}
\aulainfo{19}{Matrizes}{Tabelas de dados experimentais e pequenas simulações}

\begin{document}
\capa

\begin{frame}{Roteiro}
  \begin{itemize}
    \item Dados em tabela: linhas e colunas (Cap.~5.10)
    \item Listas de listas (5.11)
    \item Percorrendo e construindo matrizes
    \item O erro clássico: \texttt{[[0]*3]*4}
    \item Aplicações: tabelas experimentais e uma pequena simulação
  \end{itemize}
\end{frame}

\begin{frame}{Matrizes e tabelas de dados}
  \begin{columns}[T]
    \column{0.5\textwidth}
    Dados experimentais costumam vir em \textbf{tabelas}:
    \vspace{1ex}
    \centering\small
    \begin{tabular}{c|ccc}
      Ensaio & \(V\) (V) & \(I\) (A) & \(T\) (\textdegree C)\\\hline
      1 & 5,0 & 0,25 & 21,3\\
      2 & 10,0 & 0,52 & 22,1\\
      3 & 15,0 & 0,76 & 23,4\\
      4 & 20,0 & 1,01 & 25,0\\
    \end{tabular}
    \column{0.45\textwidth}
    \begin{itemize}
      \item Cada \textbf{linha} é um ensaio
      \item Cada \textbf{coluna} é uma grandeza
      \item Em matemática: matriz \(4\times3\)
      \item Em Python: uma \textbf{lista de listas}
    \end{itemize}
  \end{columns}
\end{frame}

\begin{frame}[fragile]{Lista de listas}
\begin{lstlisting}
ensaios = [
    [5.0, 0.25, 21.3],
    [10.0, 0.52, 22.1],
    [15.0, 0.76, 23.4],
    [20.0, 1.01, 25.0],
]
print(ensaios[1])         # [10.0, 0.52, 22.1]   (uma linha)
print(ensaios[1][2])      # 22.1   (linha 1, coluna 2)
print(len(ensaios))       # 4 linhas
print(len(ensaios[0]))    # 3 colunas
\end{lstlisting}
  \begin{itemize}
    \item Dois índices: \texttt{matriz[linha][coluna]}
    \item Cada linha é uma lista independente (e poderia ter tamanho diferente!)
  \end{itemize}
\end{frame}

\begin{frame}[fragile]{Percorrendo uma matriz}
\begin{lstlisting}
for i in range(len(ensaios)):
    for j in range(len(ensaios[i])):
        print(ensaios[i][j], end='\t')
    print()
\end{lstlisting}
  \textbf{Ou diretamente:}
\begin{lstlisting}
for linha in ensaios:
    for valor in linha:
        print(valor, end='\t')
    print()

# resistência (V/I) de cada ensaio
resistencias = [linha[0] / linha[1] for linha in ensaios]
# média de cada coluna
medias = [sum(l[j] for l in ensaios) / len(ensaios) for j in range(3)]
\end{lstlisting}
\end{frame}

\begin{frame}[fragile]{Construindo uma matriz}
\begin{lstlisting}
# matriz 3 x 4 de zeros, usando laços aninhados
matriz = []
for i in range(3):
    linha = []
    for j in range(4):
        linha.append(0)
    matriz.append(linha)

# o mesmo com compreensão aninhada
matriz = [[0 for j in range(4)] for i in range(3)]
print(matriz)     # [[0, 0, 0, 0], [0, 0, 0, 0], [0, 0, 0, 0]]
\end{lstlisting}
  \begin{itemize}
    \item Exercício 5.27: criar uma matriz bidimensional usando dois comandos \texttt{for}
    \item Exercício 5.28: matriz \alert{tridimensional} com compreensão (lista de listas de listas)
  \end{itemize}
\end{frame}

\begin{frame}[fragile]{O erro clássico}
\begin{lstlisting}
matriz = [[0] * 3] * 4
print(matriz)          # [[0, 0, 0], [0, 0, 0], [0, 0, 0], [0, 0, 0]]
matriz[1][0] = 1
print(matriz)
\end{lstlisting}
\begin{lstlisting}[style=saida]
[[1, 0, 0], [1, 0, 0], [1, 0, 0], [1, 0, 0]]
\end{lstlisting}
  \begin{alertblock}{O que aconteceu?}
    \texttt{[x] * 4} repete a \emph{mesma referência} quatro vezes: as 4 linhas são o \textbf{mesmo objeto} (aliasing da Aula 18). Use a compreensão: \texttt{[[0]*3 for i in range(4)]}.
  \end{alertblock}
  \fonte{code/Cap5/prog.5.b\_testematriz.py}
\end{frame}

\begin{frame}[fragile]{Operações com matrizes}
\begin{lstlisting}
def soma(a, b):
    """Soma duas matrizes de mesmas dimensões."""
    return [[a[i][j] + b[i][j] for j in range(len(a[0]))]
            for i in range(len(a))]

def transposta(a):
    return [[a[i][j] for i in range(len(a))]
            for j in range(len(a[0]))]

A = [[1, 2, 3], [4, 5, 6]]
print(transposta(A))    # [[1, 4], [2, 5], [3, 6]]
\end{lstlisting}
  \small Para álgebra linear séria (produto de matrizes, inversa, autovalores), usaremos \textbf{NumPy} na próxima aula.
\end{frame}

\begin{frame}{Pequena simulação: um tabuleiro}
  \begin{columns}[T]
    \column{0.55\textwidth}
    \begin{itemize}
      \item Matrizes representam \textbf{estados}: tabuleiros, mapas, grades de células
      \item Exemplo: o ``jogo da velha'' como matriz \(3\times3\)
      \item Valores: \texttt{' '}, \texttt{'X'}, \texttt{'O'}
    \end{itemize}
    \column{0.4\textwidth}
    \begin{block}{Funções do programa}
      \small
      \texttt{novo\_tabuleiro()}\\
      \texttt{imprime(tab)}\\
      \texttt{joga(tab, i, j, sim)}\\
      \texttt{venceu(tab, sim)}
    \end{block}
  \end{columns}
\end{frame}

\begin{frame}[fragile]{Jogo da velha: estrutura}
\begin{lstlisting}
def novo_tabuleiro():
    return [[' '] * 3 for _ in range(3)]

def imprime(tab):
    for linha in tab:
        print(' | '.join(linha))
        print('-' * 9)

def venceu(tab, s):
    linhas = any(all(c == s for c in l) for l in tab)
    colunas = any(all(tab[i][j] == s for i in range(3)) for j in range(3))
    diag = all(tab[i][i] == s for i in range(3)) or \
           all(tab[i][2 - i] == s for i in range(3))
    return linhas or colunas or diag
\end{lstlisting}
  \small Cada função faz uma coisa; o jogo completo é um desafio para o laboratório.
\end{frame}

\begin{frame}[fragile]{Simulação: espalhando calor em uma barra}
\begin{lstlisting}
# 1 linha por instante de tempo, 1 coluna por posição
n, passos = 10, 5
T = [[0.0] * n for _ in range(passos)]
T[0][0] = 100.0                       # extremidade quente
for t in range(1, passos):
    for x in range(n):
        esq = T[t-1][x-1] if x > 0 else T[t-1][x]
        dir_ = T[t-1][x+1] if x < n-1 else T[t-1][x]
        T[t][x] = (esq + T[t-1][x] + dir_) / 3
for linha in T:
    print([round(v, 1) for v in linha])
\end{lstlisting}
  \small Cada instante é calculado a partir do anterior: a matriz guarda o \emph{histórico} da simulação. (Esse tipo de problema é muito mais simples em NumPy.)
\end{frame}

\begin{frame}{Atividade}
  \begin{enumerate}
    \item Armazene a tabela de ensaios (tensão, corrente, temperatura) em uma matriz. Calcule a resistência de cada ensaio e a resistência média.
    \item Escreva \texttt{produto(a, b)}: produto matricial de duas matrizes quadradas \(n\times n\).
    \item Escreva uma função que verifica se uma matriz quadrada é \textbf{simétrica}.
    \item Complete o jogo da velha (entrada das jogadas e detecção de empate).
  \end{enumerate}
  \vspace{1ex}
  \small Desafio opcional (livro, Cap.~5): o \textbf{Jogo da Vida} de Conway é uma simulação em matriz; veja \texttt{code/Cap5/prog.5.c\_lifegame.py}.
\end{frame}

\exercicios{5}{%
  \begin{itemize}\small
    \item Listas de listas: \textbf{5.27} (matriz com \texttt{for}), \textbf{5.28} (matriz tridimensional)
    \item Revise também \textbf{5.5} e \textbf{5.9} (listas e fatias)
  \end{itemize}}

\resumo{
  \item Matriz = lista de listas; acesso \texttt{m[i][j]}
  \item Use compreensão aninhada para criar matrizes: \texttt{[[0]*c for \_ in range(l)]}
  \item Nunca use \texttt{[[0]*c]*l}: as linhas ficam compartilhadas
  \item Próxima aula: dicionários e tuplas
}

\end{document}
